\chapter{Fields, boundaries and comparison}
\label{mht:ch:fields-boundaries}

A bulk field can act on a brane without being determined by that
action. Recovering the field theory requires enough boundary
observables to detect its differential, products and higher operations.
Traces of matrices give concrete tests of such an action. Boundary
conditions and quantum contractions determine what those tests retain.
The open--closed problem is to construct the action and determine the
information it recovers.

The polar model imposes a first-jet boundary condition. On a marked
sphere, a function is required to have zero first holomorphic
derivative. The resulting Dolbeault complex has an extra cohomology
class, detected by a meromorphic one-form with a double pole.

The physical matrix fields reduce classically to a surface connection
and an adjoint scalar. The map extracting that scalar is nonlinear. At the
quantum level, eliminating the other fields also contracts observables
with one another.

Hamiltonian fields on a holomorphic symplectic surface provide a second
model. Their matrix probes retain commutator equations, relations among
those equations and central rank. Comparing the two models, or comparing
either with a closed string field theory, requires a map between their
specified fields and operations.

\section{Mixed differential forms and local observables}

Let $U$ be an oriented real surface and $Z$ a complex manifold. Mixed
forms on $U\times Z$ are sections of
\[
 \mathcal A_{U,Z}
   =\Omega^\bullet(U)\widehat\otimes\Omega^{0,\bullet}(Z).
\]
The tensor product is the completed projective tensor product of the
usual complex Fréchet spaces of smooth forms. On product opens it gives
smooth forms in both sets of variables. Compactly supported smooth
spaces carry their usual inductive-limit topologies.

On a noncompact space, integrals below use compact support.
For homogeneous forms $\xi,\xi'$ on $U$ and $\eta,\eta'$ on $Z$,
the total differential and multiplication have the tensor signs
\begin{align}
 D(\xi\otimes\eta)
   &=d_U\xi\otimes\eta+(-1)^{|\xi|}\xi\otimes\bar\partial\eta,
       \label{mht:eq:mixed-differential}\\
 (\xi\otimes\eta)(\xi'\otimes\eta')
   &=(-1)^{|\eta||\xi'|}(\xi\wedge\xi')\otimes(\eta\wedge\eta').
       \label{mht:eq:mixed-product}
\end{align}
Degrees are cohomological, and $V[k]^j=V^{j+k}$. Thus the shift
$[1]$ lowers the degree of a field vector by one. Coordinates on a
graded field space have the opposite degrees.

The adjective \emph{mixed} refers to these two differentials. For a
real vector field $v$ on $U$, contraction gives
$[d_U,\iota_v]=\mathcal L_v$. Antiholomorphic contraction gives the
corresponding identity for $\bar\partial$ on a complex coordinate
chart. Extending these identities to an interacting theory requires
compatibility with its interaction and boundary conditions. At the
quantum level it also requires compatibility with renormalization.

A classical Batalin--Vilkovisky (BV) model specifies a complex of fields,
a symplectic pairing of degree $-1$, and an action whose Hamiltonian differential
squares to zero. An infinite-dimensional pairing has a domain: a
continuous functional need not have a smooth Hamiltonian vector field.
The admissible observables and their coefficient topology are therefore
part of the model.

Write $Q$ for the free differential and $\{-,-\}$ for the degree-one
BV bracket. Let $\hbar$ be a formal parameter of degree zero.
When a regulated BV Laplacian $\Delta$ is defined, the quantum master
equation for an interaction $I$ takes the form
\[
 QI+\frac12\{I,I\}+\hbar\Delta I=0.
\]
The equation is asserted only on a domain where its three terms
are defined.

Local observables on an open set $V$ form a cochain complex
$\operatorname{Obs}(V)$. A unital prefactorization structure supplies chain
maps
\[
 \operatorname{Obs}(V_1)\otimes\cdots\otimes
 \operatorname{Obs}(V_r)\longrightarrow\operatorname{Obs}(V)
\]
for pairwise disjoint $V_i\subset V$, compatible with nested insertions
and with permutations carrying the Koszul signs. Empty insertion gives
the unit. Factorization also requires descent for Weiss covers:
every finite subset of the open lies in some member of the cover.

Local constancy in the real directions means that the maps
associated to inclusions of sufficiently small real disks are
quasi-isomorphisms, with the holomorphic open fixed. It is a property
of these observable maps. A formula for the differential on fields
alone does not establish it.

Support fixes the variance of the maps. A smooth field on $V$ restricts
to a smaller open. A compactly supported distributional coefficient of
an observable extends by zero to a larger open. These maps have
opposite directions.

The continuous dual of all smooth test functions consists of compactly
supported distributions. The dual of compactly supported smooth tests
contains distributions of arbitrary support. Each current construction specifies which
test space it uses.

\section{The polar field theory}

Take $C=\mathbb P^1$, with marked point $o=0$ in the coordinate $z$,
and fix $\kappa\in\mathbb C^\times$. For $f\in C^\infty(C)$, put
\[
 \alpha=\kappa z^{-2}dz,\qquad j(f)=\partial_zf(o),
 \qquad
 E^0=\ker j,\quad E^1=\Omega^{0,1}(C),\quad d_E=\bar\partial.
\]
All functions and forms here are globally smooth, and $E^k=0$ outside
degrees zero and one. The condition $j=0$
is imposed on degree-zero elements. It is preserved by multiplication,
so $E$ is a commutative differential graded algebra.

Its degree-minus-one
trace is the circular principal value
\[
 \tau_E(\eta)=\frac1{2\pi i}
       \lim_{r\downarrow0}\int_{C\setminus\{|z|<r\}}\alpha\wedge\eta,
 \qquad \tau_E(E^0)=0.
\]
The marked coordinate and the shrinking circles are part of the
prescription. For every smooth function $f$, polar Stokes gives
\begin{equation}\label{mht:eq:entrance-polar-stokes}
 \tau_E(\bar\partial f)=\kappa j(f).
\end{equation}
The trace is therefore closed on $E$. Chapter~\ref{mht:ch:polar-neumann}
proves its continuity and computes the resulting cohomology.

For a positive integer $N$, matrix multiplication and the ordinary
trace give the classical field space
\begin{equation}\label{mht:eq:entrance-polar-fields}
 \mathcal F_E
   =\bigl(\Omega^\bullet(U)\widehat\otimes E
                         \otimes\operatorname{Mat}_N(\mathbb C)\bigr)[1].
\end{equation}
Integration over $U$, followed by $\tau_E\operatorname{tr}$, has degree
$-3$ before the shift. The induced BV pairing has degree $-1$.

Let $\mathsf A$ denote the full graded field, including its ghost and
antifield components. For a closed surface, or compact support in its
interior, the cyclic action is
\begin{equation}\label{mht:eq:entrance-polar-action}
 S_E(\mathsf A)=\int_U\tau_E\operatorname{tr}
       \left(\frac12\mathsf A D\mathsf A+\frac13\mathsf A^3\right),
 \qquad D=d_U+d_E.
\end{equation}
The differential graded algebra identities give its classical gauge
complex. This is the partial-connection form of four-dimensional
Chern--Simons theory, with the polar trace and boundary domain just
specified.\footnote{K.~Costello, E.~Witten and M.~Yamazaki,
\emph{Gauge Theory and Integrability, I},
\href{https://arxiv.org/abs/1709.09993}{arXiv:1709.09993},
Section~3.1, especially equations~(3.1)--(3.4). The action here uses
the normalization $\frac12\mathsf A D\mathsf A+\frac13\mathsf A^3$
and the displayed principal-value trace.}

A physical field is $\mathsf A=b+a$, where $b$ is a surface one-form
with coefficients in $E^0$ and $a$ is a curve $(0,1)$-form. Gauge
transformations take values in invertible matrices over $E^0$ and act by
\[
 \mathsf A^g=g^{-1}\mathsf Ag+g^{-1}Dg.
\]
Unless stated otherwise, gauges and fields are considered formally
near the zero field and the trivial curve bundle. The analytic
normal form is asserted on its separately specified small Sobolev
neighborhood.

The residual curve algebra is
\[
 A_\kappa=\mathbb C\oplus\mathbb C\varepsilon,
 \qquad |\varepsilon|=1,\quad\varepsilon^2=0,\quad
 d_A=0,\quad \tau_A(1)=0,\quad \tau_A(\varepsilon)=\kappa.
\]
It is represented in $E$ by
\[
 f_0=\frac{z}{1+|z|^2},\qquad
 \eta_0=\bar\partial f_0,
 \qquad 1\longmapsto1,\quad\varepsilon\longmapsto\eta_0.
\]
Although $\eta_0$ is exact among unrestricted smooth forms,
\eqref{mht:eq:entrance-polar-stokes} excludes a Neumann primitive.
The inclusion is a trace-preserving quasi-isomorphism.

The residual matrix of a curve connection is extracted from its
normalized frame:
\[
 \bar\partial H=-aH,\qquad H(o)=I_N,\qquad
 \mathcal J(a)=-j(H).
\]
Let $G$ be the right inverse of $\bar\partial$ normalized by
$(G\eta)(o)=0$. Formally,
\begin{equation}\label{mht:eq:entrance-nonlinear-residual}
 H=I_N-Ga+G(aGa)-\cdots,\qquad
\mathcal J(a)=j(Ga)-jG(aGa)+\cdots.
\end{equation}

The order of the matrix factors is retained. The first term is the
linear cohomology projection. The other terms record how the
connection changes that projection. Chapter~\ref{mht:ch:nonlinear-neumann}
constructs the frame, its equivariance and explicit nonzero nonlinear
trace invariants.

For a surface connection $b\in\Omega^1(U,\mathfrak{gl}_N)$ and
$J\in C^\infty(U,\mathfrak{gl}_N)$, the residual field is
$b+\eta_0J$. The same physical density on a surface with boundary gives
\begin{equation}\label{mht:eq:entrance-bf-action}
 S_E(b+\eta_0J)
   =\kappa\int_U\operatorname{tr}(JF_b)
        -\frac\kappa2\int_{\partial U}\operatorname{tr}(bJ),
 \qquad F_b=d_Ub+b\wedge b.
\end{equation}
The boundary term contributes to the variational problem. A boundary
polarization must also be chosen.

The full
shifted residual complex also contains ghosts and antifields. At a
constant matrix $J$, its curve differential is the two-term adjoint
complex $[\mathfrak{gl}_N\xrightarrow{\operatorname{ad}_J}
\mathfrak{gl}_N\varepsilon]$. Retaining this complex preserves the
stabilizers when eigenvalues coincide. Chapter~\ref{mht:ch:neumann-bv}
proves the full reduction and the two boundary transgressions, at the
pole and at the boundary of $U$.

The coefficient ring matters in this reduction. Formal deformation
parameters admit recursive inversion, and sufficiently small analytic
parameters admit the corresponding convergent inverse. A polynomial
parameter does not permit either step. The same chapter exhibits an
infinite-dimensional cokernel for the polynomial family while retaining
the formal and analytic equivalences. Thus the residual model comes
with its parameter domain.

\section{Hamiltonian fields and matrix probes}

The symplectic model has a different holomorphic factor. Let $W$ be a
complex surface with holomorphic symplectic form. On a Darboux chart
write this form as $dq\wedge dp$ and use
\[
 \{f,g\}=f_qg_p-f_pg_q,\qquad X_f=\{f,-\}.
\]
The spacetime $U\times W$ has real dimension six, whereas $U\times C$
in \eqref{mht:eq:entrance-polar-fields} has real dimension four.
A passage between them requires a reduction with its fields,
integration domain and boundary conditions.

At a fixed point of $W$, put $\mathfrak m=(q,p)$ in
$\mathbb C[[q,p]]$. Formal Hamiltonian vector fields are represented
by $\mathbb C[[q,p]]/\mathbb C1$. Those fixing the point are represented
by $\mathfrak m^2$. The latter have finite Lie quotients
\[
 \mathfrak h_d=\mathfrak m^2/\mathfrak m^{d+1},\qquad d\geq2,
\]
because the Poisson bracket lowers polynomial degree by two.
Translations do not preserve these quotients:
$\{q,p^{d+1}\}=(d+1)p^d$. The pointed and moving-point problems
therefore require different coefficient data.

An internal formal Hamiltonian is also distinct from a function of
the physical holomorphic coordinates. If $z=(z_1,z_2)$ are physical
coordinates on a chart $V$ and $w=(w_1,w_2)$ independent formal variables,
use the complex
\[
 \mathcal O(V)[[w_1,w_2]]\otimes\Lambda(\vartheta_1,\vartheta_2),
 \qquad |\vartheta_i|=1.
\]
Its horizontal jet differential is
\[
 \nabla=\sum_{i=1}^2\vartheta_i
              (\partial_{z_i}-\partial_{w_i}).
\]
Its horizontal inclusion sends $f(z)$ to the formal Taylor series
$f(z+w)$. Evaluation at $w=0$ on arbitrary coefficients is not a
Poisson map: it sends $\{w_1,w_2\}_w$ to $1$ and sends both inputs
to zero. The horizontal complex supplies the additional constraint
needed for the geometric Hamiltonian comparison.

The matrix probe of a commutative plane is the differential graded
algebra
\begin{equation}\label{mht:eq:entrance-matrix-probe}
 R_N=\mathbb C[X_{ij},Y_{ij}]\otimes\Lambda(\Psi_{ij}),\qquad
 |\Psi_{ij}|=-1,\quad Q\Psi=[X,Y],\quad QX=QY=0.
\end{equation}
The Chevalley--Eilenberg cochains for simultaneous conjugation additionally
record infinitesimal gauge symmetry and stabilizers.
Even before that step, the derived relation matters:
$\operatorname{Tr}\Psi$ is a closed class that is not exact. Evaluation
at $X=Y=0$ sends it to a nonzero element of the exterior algebra
with zero differential, proving the assertion. A reduction that
removes this class changes the probe.

Central constants are detected independently. With
$\{X_{ij},Y_{kl}\}=\delta_{il}\delta_{jk}$, direct summation gives
\begin{equation}\label{mht:eq:entrance-rank-bracket}
 \{\operatorname{Tr}X,\operatorname{Tr}Y\}=N.
\end{equation}
The classes of $q,p$ commute in the Lie algebra of Hamiltonians modulo
constants. Their unreduced matrix traces have the nonzero bracket
\eqref{mht:eq:entrance-rank-bracket}. A Hamiltonian-function comparison
must retain this central datum or construct its reduction. The
induced derivations, which annihilate constants, ask a different
question.

Framed probes impose a moment equation with an additional scalar
parameter. At rank two, a cyclic matrix basis gives regular coordinates
even when eigenvalues collide. Chapter~\ref{mht:ch:matrix-branes}
constructs an integral quantum reduction on this chart and a regular
Lie map from polynomial Hamiltonians.

Two different global obstructions remain. One fixes four low charges. The other allows
every formal correction but prescribes the full trace map at zero
framing. A comparison at an isolated nonzero framing level can
lie outside both classes. Its nonlinear moments and its transitions
between cyclic charts must still be constructed.

\section{The bulk action and the boundary center}

A boundary theory includes objects describing boundary conditions and
cochain complexes of operators changing one condition into another.
Their compositions are part of the data. For an associative
differential graded algebra $B$, the corresponding algebraic center is
the Hochschild cochain complex
\[
 C^\bullet(B,B)=\operatorname{RHom}_{B\otimes B^{\mathrm{op}}}(B,B).
\]

It records coherent operations on $B$ and its modules. A geometric
boundary theory requires the analogous operations to respect the
positions and support of insertions. Choosing its category of
factorization bimodules and its diagonal precedes taking derived
endomorphisms. That choice is not determined by the trace of a single
matrix probe.

Even the polynomial coordinate comparison has two stages. Regrading
the Chevalley--Eilenberg coordinates of a shifted cotangent Lie algebra
gives polyvectors with their linear Poisson differential. Ordinary
Hochschild--Kostant--Rosenberg antisymmetrization has zero differential
on its polyvector source. It therefore supplies the Hochschild
comparison only at zero Poisson structure. A nonzero Poisson tensor
requires a supplied formality map, twisted by that tensor, and the
corresponding deformed associative product.
Chapter~\ref{mht:ch:centres-actions} proves the coordinate identity
and gives the explicit failure of the untwisted comparison.

The bulk is specified independently by its fields, pairing, action
and boundary domain. Its action on a boundary is a family of maps
from bulk observables in a collar to coherent boundary operations.
The maps must commute with the differential and preserve disjoint
products. Moving and colliding insertions impose further homotopies.

For example, let $\mathfrak l$ be a Lie algebra in degree zero and
$(\mathfrak a,d)$ a differential graded Lie algebra.
Let $F_1:\mathfrak l\to\mathfrak a^0$ take values in closed elements.
A binary correction is a map
$F_2:\Lambda^2\mathfrak l\to\mathfrak a^{-1}$ required to solve
\[
 dF_2(f,g)=F_1([f,g])-[F_1(f),F_1(g)].
\]
The next equation tests the compatibility of these binary corrections.
An action on cohomology alone leaves these equations undecided.

A dynamical coupling must also determine the source in the bulk
equations. If $i:Y\hookrightarrow M$ is a defect inclusion, the
transpose of restriction sends a defect density $\mu$ to a current
supported on $Y$. For a real normal coordinate $r$, oriented before
the coordinates of $Y$, this pushforward has the form
$i_*\mu=\delta(r)\,dr\,\mu$. Coupling a Hamiltonian bulk field to a
cyclic action on defect fields produces both their gauge action and
this transposed source. A product of singular currents at several
insertions requires its own extension across their collisions.

The geometric Hamiltonian model uses a six-dimensional bulk and a
five-dimensional defect of this kind. Its cotangent field admits
normal distributions, while the field multiplying it is smooth.
This domain permits the classical action without multiplying two
normal distributions. A closed theory of polyvector fields, such as
Bershadsky--Cecotti--Ooguri--Vafa (BCOV) theory, starts with a different
complex and its own divergence, descendants and background.
Identifying its sourced observables with Hamiltonian or polar boundary
operations requires a map of those complexes, preserving the source
terms and the cyclic operations. The general comparison remains open.

\section{Quantum reduction and prepared states}

Classical reduction supplies a residual field and the homotopies of
its gauge complex. Quantum elimination additionally supplies
contractions. For a normalized Gaussian functional on one variable $x$,
with mean $m$ and covariance $\hbar C$,
\[
 \langle x^2\rangle=m^2+\hbar C.
\]
Substitution of $m$ alone misses the second term. For several
observables the corresponding corrections are their connected
cumulants. The normalized BV pushforward and its composition law in
Chapter~\ref{mht:ch:quantum-elimination} retain these terms.

Infinite-dimensional elimination requires more than the finite
identity. The paired native heat tensors give continuous contractions
at positive scale, including pole jets. Their generator is nonlocal
on the curve. Extending coincident products to zero scale therefore
requires estimates and counterterms on the stated support domain.
The free identities of Chapter~\ref{mht:ch:heat-contacts} do not supply
the interacting quantum master equation.

That chapter also computes the full supported normal-jet complex.
Antiholomorphic derivative jets have explicit Dolbeault primitives.
Holomorphic residue jets survive in cohomology. Boundary faces and
their intersections require the relative differential. Admissibility
of these primitives for the native cyclic field theory is a further
condition.

For the Gaussian sector, the staggered inverse converges uniformly on
compact sets of matrix backgrounds. Its finite background derivatives
converge on bounded smooth tests. Chapter~\ref{mht:ch:quantum-elimination}
uses these estimates to obtain smooth Ward identities and the
regularized determinant. An inserted source detector remains nonzero
even when the uninserted regularized determinant equals one.

A state is further data: a functional on observables compatible with
their differential and the chosen boundary preparation.
The unnormalized functional evaluated on the unit gives the partition
function. Normalizing that value leaves the other observable values
to be determined.

Chapter~\ref{mht:ch:states-sewing} constructs compact
gauge and Gaussian preparations, keeps their normalization lines,
and proves their sewing formulas. Its interacting scalar preparation
has a fixed positive smoothing parameter. A Hermite unitary identifies
it with an independently defined scalar Fock interaction, including
the full form domains, semigroups, traces and normalized positive
states. The finite cutoff semigroups converge in trace norm.

The remaining physical comparison must identify the state produced by
the native field theory with an independently prepared residual state
on common observables. It must also intertwine cutting and sewing.
This requirement connects the boundary algebra to the quantum field
theory: the algebra specifies which operations can be composed, and
the preparation specifies their amplitudes.
